\uuid{Xijj}
\exo7id{4564}
\titre{exo7 4564}
\auteur{quercia}
\organisation{exo7}
\datecreate{2010-03-14}
\isIndication{false}
\isCorrection{false}
\chapitre{Série entière}
\sousChapitre{Rayon de convergence}
\cpge{mpsi-mp2i.procedes-sommatoires-discrets.a, pcsi.series-numeriques.a, ptsi.series-numeriques.a, mpsi-mp2i.procedes-sommatoires-discrets.c, pcsi.series-numeriques.c, ptsi.series-numeriques.c, bcpst1.suites-reelles, mp-mpi.series-entieres.a, pc.series-entieres.a, psi.series-entieres.a, pt.series-entieres.a, mp-mpi.series-entieres.b, pc.series-entieres.d, psi.series-entieres.d, pt.series-entieres.d, mp-mpi.series-entieres.c, pc.series-entieres.b, psi.series-entieres.b, pt.series-entieres.b, bcpst2.series-reelles}
\module{Analyse}
\niveau{L2}
\difficulte{}

\contenu{
\texte{
Dire si les affirmations suivantes sont vraies ou fausses. En donner
une démonstration ou un contre-exemple.
}
\begin{enumerate}
    \item \question{Les séries $\sum a_nz^n$ et $\sum (-1)^na_nz^n$ ont même rayon de
    convergence.}
    \item \question{Les séries $\sum a_nz^n$ et $\sum (-1)^na_nz^n$ ont même domaine de
    convergence.}
    \item \question{Si la série $\sum a_nz^n$ a un rayon de convergence infini, alors elle
    converge uniformément sur $\R$.}
    \item \question{Si $\sum a_nx^n$ a un rayon de convergence fini $R > 0$, alors sa somme
    admet une limite infinie en $(-R)^+$ ou en~$R^-$.}
    \item \question{Si $f(x) = \sum a_nx^n$ a un rayon de convergence infini et si les $a_n$
    sont strictement positifs, alors pour tout entier~$p$,
    $\frac{f(x)}{x^p} \to +\infty$ lorsque $x\to+\infty$.}
\end{enumerate}
}
